\section{The initial enclosure, buried boundary, and flow notation}
\label{sec:rpi:flow-localization}

We construct the conformal flow for $n\ge8$, following Bray and
Bi--Zhu \cite{Bray2001,BiZhu2026v2}. The flow keeps the full frontier
area fixed while changing the exterior metric. We first specify its
ambient manifold and filled side, including the connectedness needed
for the harmonic potentials.

Extend the original smooth metric a short distance behind its compact
frontier, and truncate there at a smooth compact boundary $B$. Denote
the resulting smooth ambient exterior by $(M,g_0)$ and its original
filled side by $J$. The set $J$ contains a fixed collar of $B$.
Smooth boundary charts, the compact core, and the complete
asymptotically flat tail make $M$ complete with $B$ included. The
metric on the buried side need not have nonnegative scalar curvature.

Use the canonical closed representative of $J$, whose frontier is the
full perimeter support. Its open complement is connected. Indeed,
there is one component containing the coordinate end, and any other
component $U$ is bounded with boundary in the original frontier.
That frontier has finite $(n-1)$-measure and a singular part of zero
$(n-1)$-measure, so the finite-boundary-measure criterion gives finite
perimeter for $U$. On regular boundary charts its normal is opposite
to the filled-side normal. Filling $U$ therefore gives
\[
 P_{g_0}(J\cup U)=P_{g_0}(J)-P_{g_0}(U).
\]
The nonempty bounded open set $U$ has positive perimeter, contradicting
outer minimization of $J$. This proves connectedness of the open end
exterior. The perimeter facts used here are local BV facts
\cite{AmbrosioFuscoPallara2000,Maggi2012}.

Local perimeter minimality of the original frontier and its compactness
give one radius for all sufficiently small local comparisons. Shrink
that radius so the variation balls also avoid the buried collar.
Balls missing the frontier have zero local perimeter. We may replace
$J$ by its outermost minimizing enclosure $K_0$: outer minimization
preserves its area, and the lattice comparison below preserves its
local minimality. Write $E_0=M\setminus K_0$. The replacement has
connected open exterior by the same filling argument. Scalar
nonnegativity holds on $E_0$, and later it will be used on
$E_t\subset E_0$. The auxiliary physical boundary remains buried
throughout the flow.
With $k=n-1$ and $d=n-2$, the flow constructed below has $E_t\subset E_0$ and
defining equations
\[
 u_t=1+\int_0^t v_s\,ds,\qquad g_t=u_t^{4/d}g_0,
 \qquad \Delta_{g_0}v_t=0\ \hbox{in }E_t,
 \quad v_t=0\ \hbox{on }K_t,
 \quad v_t\longrightarrow-e^{-t}.
\]
Consequently $v_t=0$ and $u_t=1$ on the same fixed collar of $B$.
The $g_0$-capacitary potential is $1+e^t v_t$; the $g_t$-capacitary
potential, denoted below by $w_t$, is
\[
 w_t=1+v_t/u_t\quad\hbox{on }E_t,
 \qquad w_t=1\quad\hbox{on }K_t.
\]
These two potentials must be distinguished. Both are constant on that
collar and have zero physical-boundary flux. Capacity test functions
are compactly supported relative to $M$, smooth up to $B$, and equal
to one near $K_t$; they are not required to vanish on $B$. Equivalently,
one uses the Dirichlet problem on $E_t$ and extends by one to $K_t$.
Every regular level below one excludes the buried collar, so the
capacity-flux calculation has only its level boundary and its sphere
at infinity. No mean-curvature condition at $B$ enters.

The smooth reference metric remains $g_0$. The working metric $g_t$
is smooth on the open exterior, where $u_t$ is $g_0$-harmonic and
$R_{g_t}=u_t^{-4/d}R_{g_0}\geq0$. Across a singular frontier one uses
the comparison and H\"older estimates of
Lemma~\ref{lem:rpi:flow-construction}, followed by the simultaneous
one-sided regularity and relative forcing estimates of
Section~\ref{sec:rpi:regular-patches}. The perimeter-plus-volume forcing problem
uses filled-set competitors containing $K_t$; they may meet its
frontier. Corollary~\ref{cor:rpi:flow-detachment} proves detachment
including at singular points. Once
its new frontier is strictly detached, it has a neighborhood compactly
contained in this smooth exterior. Its one-sided smoothing and the
smooth mass--capacity argument therefore take place entirely outside
the original frontier and omit $B$ altogether.

\section{Construction and preserved area of the flow}
\label{sec:rpi:flow-construction}

\begin{lemma}[Construction and area of the conformal flow]
\label{lem:rpi:flow-construction}
Let $K_0$ be the initial outermost minimizing enclosure just described,
and put
\[
 E_0=M\setminus K_0,\qquad \Sigma_0=\partial K_0,
 \qquad A_0=P_{g_0}(K_0).
\]
Here and below perimeter is relative to $\operatorname{int}M$ and
counts the whole frontier. There are increasing closed filled sets
$K_t\supset K_0$, with connected end exteriors $E_t=M\setminus K_t$,
and positive continuous functions $u_t$ such that
\begin{equation}
\label{eq:rpi:flow-integral}
 \begin{gathered}
 g_t=u_t^{4/d}g_0,\qquad
 u_t=1+\int_0^t v_s\,ds,\\
 \Delta_{g_0}v_t=0\quad\hbox{in }E_t,\qquad
 v_t=0\quad\hbox{on }K_t,\qquad
 v_t\longrightarrow-e^{-t}\quad\hbox{at infinity}.
 \end{gathered}
\end{equation}
For every $t$ we have $e^{-t}\le u_t\le1$. Each $K_t$ locally
minimizes perimeter in the continuous metric $g_t$
against all variations in balls below the initial comparison radius,
including variations across $K_0$. The set $K_t$ is the
outermost $g_t$-minimizing enclosure of $K_0$ and satisfies
\begin{equation}
\label{eq:rpi:flow-area}
 P_{g_t}(K_t)=A_0.
\end{equation}
On each finite time interval the conductors have common compact
support, the factors have a common positive lower bound and spatial
H\"older modulus, and the frontiers have uniform two-phase density
and almost-minimizing estimates. In particular there are
$\alpha_T>0$, $C_T<\infty$, and $r_T>0$ such that
\begin{equation}
\label{eq:rpi:flow-almost-minimizing}
 P_{g_0}(K_t;B_r(x))
 \le P_{g_0}(L;B_r(x))+C_T r^{k+\alpha_T}
\end{equation}
whenever $0\le t\le T$, $r<r_T$, and
$L\mathbin{\triangle}K_t\Subset B_r(x)$.
\end{lemma}


The discrete method is due to Bray \cite{Bray2001}; its singular-frontier
form follows \cite[Section~2.1]{BiZhu2026v2}. We follow the order of
\cite[Lemmas~2.21 and~2.23 and Corollary~2.22]{BiZhu2026v2}:
area preservation precedes nesting of the limiting outermost hulls.
Freezing on a discrete conductor by itself does not give freezing on
a potentially larger limiting conductor.

The proof has three stages. A static weighted enclosure problem supplies
compactness and boundary estimates for a single conductor. We apply
those estimates to the discrete flow and prove that its area loss tends
to zero. Finally, we identify the minimum area of each limiting metric;
that equality gives nesting and the integral evolution.

\subsection{Weighted enclosures and their escape potentials}

For a continuous positive weight $a$, write
\[
 P_a(L;U)=\int_{U\cap\partial^*L}a\,d\mathcal H^k_{g_0}.
\]
The next lemma requires only continuity of $a$ for its density and
potential estimates. Its last assertion records the stronger comparison
available when the weight is H\"older continuous.

\begin{lemma}[A weighted enclosure and its potential]
\label{lem:rpi:static-hull-potential}
Every compact finite-perimeter obstacle has a unique outermost
minimizing enclosure for any continuous perimeter weight bounded above
and below by positive constants.

For the quantitative conclusions, let $K_0$ be a compact
finite-perimeter filled set containing a collar of $B$, locally
perimeter minimizing in $g_0$ with comparison radius $r_0>0$. Put
$A_0=P_{g_0}(K_0)$, and let $a$ be continuous with
$c\le a\le1$ and $a=1$ on $K_0$, where $c>0$.
Every $P_a$-minimizing enclosure $K$ has weighted perimeter at most $A_0$
and is contained in a ball $B_{R_*(c)}$ depending only on $c$, $A_0$,
the initial support, and the reference geometry.

Every such enclosure locally minimizes $P_a$ against unconstrained
variations in balls of radius less than $r_0$. Its canonical closed
representative has connected open end exterior and uniform two-phase
volume and perimeter density bounds. The unique escape potential
$q_K$, harmonic in that exterior, zero on $K$, and tending to one at
infinity, is continuous and has a H\"older modulus depending on the
same data. It satisfies $0<q_K<1$ in the exterior.

If in addition $\operatorname{osc}_{B_r}a\le C_a r^\alpha$ for
$0<\alpha<1$, then, for local competitors $L$,
\[
 P_{g_0}(K;B_r)\le P_{g_0}(L;B_r)+C r^{k+\alpha},
\]
where $C$ also depends on $C_a$.
\end{lemma}

\begin{proof}
\textbf{Confinement and existence.}\quad
First minimize $P_a$ among enclosures of $K_0$, either globally or
with an outer constraint $L\subset B_Q$. The competitor $K_0$ bounds
the minimizing perimeter by $A_0$.
Choose $R_0$ containing $K_0$ and the compact core, and choose
$\gamma\ge1$ such that
$\gamma^{-2}\delta\le g_0\le\gamma^2\delta$ in the coordinate
end. Put
\[
 a_-=c\gamma^{-k},\qquad a_+=\gamma^k.
\]
For almost every $R>R_0$, truncation by $B_R$ is admissible and
cancellation of the common interior perimeter gives
\begin{equation}
\label{eq:rpi:flow-truncation}
 P_\delta(L;\{|x|>R\})
 \le\frac{a_+}{a_-}\mathcal H^k_\delta(L\cap S_R).
\end{equation}
For the constrained problem this is used when $R<Q$; beyond $Q$
the tail is empty. Set $V(R)=|L\cap\{|x|>R\}|_\delta$.
Euclidean isoperimetry applied to the cut tail, including its
spherical boundary, and the coarea formula imply
\[
 V(R)^{k/n}
 \le C_{\rm iso}\left(1+\frac{a_+}{a_-}\right)[-V'(R)].
\]
At a regular cutting radius $R_1\in[R_0,R_0+1]$, a second use of
Euclidean isoperimetry gives
\[
 V(R_1)\le V_*:=
 \left[C_{\rm iso}\left(a_-^{-1}A_0+
              \mathcal H^k_\delta(S_{R_0+1})\right)\right]^{n/k}.
\]
Indeed, the original end perimeter is at most $a_-^{-1}A_0$ and
the added cutting boundary is a portion of that sphere. While
$V>0$,
\[
 \frac{d}{dR}V(R)^{1/n}
 \le-\left[nC_{\rm iso}\left(1+\frac{a_+}{a_-}\right)\right]^{-1}.
\]
Consequently the tail vanishes by the radius
\begin{equation}
\label{eq:rpi:flow-hull-radius}
 R_*(c)=R_0+1+
 nC_{\rm iso}\left(1+\frac{a_+}{a_-}\right)V_*^{1/n}.
\end{equation}
We use the canonical closed representative
$L=\supp(\chi_L\,dV_{g_0})$; zero tail volume therefore means
actual containment in $B_{R_*(c)}$. No density estimate,
connectedness assertion, or potential estimate was needed.

Minimize first in each bounded exhaustion domain $B_Q$. BV
compactness and weighted lower semicontinuity give a constrained
minimizer. The radius in \eqref{eq:rpi:flow-hull-radius} is
independent of $Q$, so a sequence with $Q\to\infty$ has a limit
on one fixed support ball. Every fixed compact competitor is
admissible eventually; the limit is thus a global minimizer.
All global minimizers have the same support bound. Choose one
of maximal background volume by compactness once more.
Submodularity and minimality show that the union and intersection
of any two minimizers are again minimizers. The maximal-volume
one must therefore contain every other minimizer, first almost
everywhere and then canonically. This is the unique outermost hull.
The same argument constructs the minimizing enclosure of any
compact finite-perimeter obstacle $J$ for any continuous weight
bounded above and below by positive constants: use $P_a(J)$ as
the competitor bound and choose $R_0$ to contain $J$.
Equality of the weight to one on $J$ is unnecessary for this
individual existence assertion.

If the prescribed initial enclosure $J$ was replaced by its
outermost minimizing hull $K_0$, the replacement also preserves
local minimality, not only area. For
$L\mathbin{\triangle}K_0\Subset U$ in an initial variation ball,
apply the local minimality of $J$ to $L\cap J$ and the enclosure
minimality of $K_0$ to $L\cup J$. Submodularity gives
\[
 \begin{aligned}
 P_{g_0}(K_0;U)
 &\le P_{g_0}(L\cup J;U)\\
 &\le P_{g_0}(L;U)+P_{g_0}(J;U)-P_{g_0}(L\cap J;U)\\
 &\le P_{g_0}(L;U).
 \end{aligned}
\]
Thus the same initial comparison radius applies to $K_0$.

\paragraph{Removal of the initial obstacle.}
Let $K$ be any minimizing enclosure. For a local competitor
$L\mathbin{\triangle}K\Subset B_r(x)$ with $r<r_0$, initial
local minimality and $a=1$ on $K_0$ give
\[
 P_a(K_0;B_r(x))
 \le P_a(L\cap K_0;B_r(x)).
\]
Global enclosure minimality applies to $L\cup K_0$.
Combining it with the local submodularity inequality yields
\begin{equation}
\label{eq:rpi:flow-local-weighted-minimality}
 P_a(K;B_r(x))
 \le P_a(L;B_r(x)).
\end{equation}
Thus $K$ is a local background quasiminimizer with factor $Q=c^{-1}$.
The comparison includes balls crossing the initial obstacle and uses
its local minimality against removals. Cancelling unchanged perimeter
also localizes the comparison to smaller variation domains.

\paragraph{Two-phase density before boundary regularity.}
Choose uniform small background coordinate balls away from a
smaller buried collar. Smoothness on the compact core and the
differentiated asymptotic-flatness bounds on the end give uniform
local absolute and relative isoperimetric constants, sphere-area
bounds, and scaled $C^{1,\alpha_0}$ bounds for the divergence-form
coefficients of $\Delta_{g_0}$, for some $\alpha_0\in(0,1)$.
No uniform bounds for derivatives of all orders are required.
Shrink their
radius below $r_0$. Balls missing
$\partial K_0$ have zero initial local perimeter and cause no
exception. We may use geodesic balls in the following slicing
argument by localization of the exact weighted comparison.

Let $K$ be one of the hulls and
$x\in\supp|D\chi_K|$. Put
\[
 m(r)=\Vol_{g_0}(K\cap B_r(x)),\qquad
 e(r)=\Vol_{g_0}(B_r(x)\setminus K).
\]
Both numbers are positive for every $r>0$, since otherwise the
perimeter measure would vanish near $x$. Delete or fill $B_r(x)$,
first comparing in a slightly larger ball, and then decrease the
comparison radius. BV slicing and quasiminimality give, for
almost every $r$,
\[
 P_{g_0}(K;B_r(x))\le Qm'(r),\qquad
 P_{g_0}(K;B_r(x))\le Qe'(r).
\]
If $I$ is the local absolute isoperimetric constant, the perimeter
bound and isoperimetric inequality give
\[
 m(r)^{k/n}\le I\bigl(P_{g_0}(K;B_r(x))+m'(r)\bigr)
 \le I(Q+1)m'(r).
\]
Integrating the resulting inequality for $m^{1/n}$ from its zero
limit at radius zero, and doing the same for the exterior phase,
gives
\[
 \min\{m(r),e(r)\}\ge[nI(Q+1)]^{-n}r^n.
\]
Also $m'(r)+e'(r)=\Area_{g_0}(\partial B_r(x))\le Sr^k$,
so adding the two comparison inequalities gives
$P_{g_0}(K;B_r(x))\le QS r^k/2$. Perimeter-measure
monotonicity extends the upper bound to every small radius;
relative isoperimetry and the two phase volumes give the lower
bound. Thus, with constants depending only on the stated data,
\begin{equation}
\label{eq:rpi:flow-density-input}
 \begin{gathered}
 \min\{\Vol_{g_0}(K\cap B_r(x)),
          \Vol_{g_0}(B_r(x)\setminus K)\}\ge c_* r^n,\\
 c_* r^k\le P_{g_0}(K;B_r(x))\le C_* r^k
 \qquad(x\in\partial K,\ 0<r<r_*).
 \end{gathered}
\end{equation}
For the canonical representative $K=\supp(\chi_K\,dV_{g_0})$,
these estimates identify $\partial K$ with
$\supp|D\chi_K|$. They hold at every frontier point, not only
almost every reduced-boundary point. Null-set additions or
artificial isolated points are not part of the representative.

A bounded complementary component of a minimizing hull can be
filled, removing its strictly positive perimeter. This contradicts
minimality; equivalently one may use the finite-perimeter
component decomposition. The canonical representatives have no
measure-zero cracks producing artificial interfaces. Since the
manifold has only one end and the conductor is compact, the
remaining exterior is its connected end component.

\paragraph{The weak capacitary problem.}
The harmonic function is constructed before any continuity at the
frontier is assumed. For an individual compact conductor $K$, let
$\mathcal D^{1,2}(M)$ be the completion, in the gradient norm, of
smooth functions compactly supported relative to $M$. Such
functions may be nonzero on the physical boundary. The
asymptotically flat Sobolev inequality, for $n\ge3$, embeds this
space in $L^{2n/d}$. Minimize
\[
 \int_M|\nabla W|_{g_0}^2\,dV_{g_0}
\]
over its closed convex subset consisting of functions whose
quasi-continuous representative is at least one quasi-everywhere
on $K$. Here quasi-everywhere means outside a set of Newtonian
capacity zero. A smooth function equal to one near $K$ makes the
class nonempty. The Hilbert-space direct method gives a unique
minimizer $W_K$; truncation gives $0\le W_K\le1$ and $W_K=1$
quasi-everywhere on $K$.

Variations of both signs supported in $M\setminus K$ prove
harmonicity there. Nonnegative variations supported away from
the physical boundary remain admissible, and give
\[
 \int_M\langle\nabla W_K,\nabla\phi\rangle_{g_0}\,dV_{g_0}
 \ge0\qquad(\phi\ge0).
\]
Consequently the escape potential $q_K=1-W_K$, extended by zero
on $K$, is locally $H^1$, satisfies $0\le q_K\le1$, and obeys
\begin{equation}
\label{eq:rpi:flow-weak-potential}
 \Delta_{g_0}q_K\ge0\quad\hbox{weakly away from the physical boundary},
 \qquad \Delta_{g_0}q_K=0\quad\hbox{in }M\setminus K.
\end{equation}
Sobolev integrability of $W_K$ and interior estimates on end
annuli imply $W_K\to0$ uniformly at infinity. A radial
supersolution $r^{-\beta}$, $0<\beta<d$, and interior gradient
estimates give, after increasing a radius depending on $K$,
\[
 W_K=O(r^{-\beta}),\qquad |\nabla W_K|=O(r^{-1-\beta}).
\]
No uniform support radius is assumed in this individual estimate.
Thus $q_K\to1$ at infinity. The homogeneous space is important:
$W_K$ need not belong to unweighted $L^2$ in dimensions three
and four. Since $K$ contains a collar of the physical boundary,
$W_K=1$ and $q_K=0$ there, and the boundary contributes no test
or flux term to this exterior problem.

\paragraph{The exterior-indicator boundary contraction.}
Fix $x\in\partial K$ and work in uniform background coordinates,
using Euclidean coordinate radii in this paragraph. By
\eqref{eq:rpi:flow-density-input} and metric comparability,
$|K\cap B_r(x)|\ge d_* r^n$. Choose $\kappa\in(0,1/2)$ so
small that $|B_{\kappa r}(x)|\le d_* r^n/2$. Coarea supplies a
radius $\rho\in[\kappa r,r]$ at which the Sobolev traces exist
and
\[
 \mathcal H^k_\delta(K\cap\partial B_\rho(x))\ge d'_* r^k.
\]
At this radius the zero extension of $q_K$ has zero trace almost
everywhere on the conductor portion of the sphere. Set
\[
 \chi=\mathbf1_{(M\setminus K)\cap\partial B_\rho(x)},
 \qquad M(r)=\operatorname*{ess\,sup}_{B_r(x)}q_K.
\]
Let $H$ be the $g_0$-harmonic replacement in the full ball with
boundary trace $q_K$. Weak subharmonicity in
\eqref{eq:rpi:flow-weak-potential} and the weak maximum principle
give $q_K\le H$. If $P_\rho$ is the Poisson kernel, then
\[
 H(y)\le M(r)\int_{\partial B_\rho(x)}
                 P_\rho(y,\xi)\chi(\xi)\,dS_\delta(\xi).
\]
The datum is the exterior indicator $\chi$, not $1-\chi$:
only the former dominates the possibly positive trace of $q_K$.
For the fixed smooth background operator, scaled smooth-ball
Green-function estimates give
\[
 P_\rho(y,\xi)\ge c_P\rho^{1-n}
 \quad(y\in B_{\rho/2}(x),\ \xi\in\partial B_\rho(x)).
\]
The constant is uniform under the local ellipticity and
coefficient bounds already fixed; the kernel integrates to one.
Its missing mass on the conductor portion of the sphere is
therefore at least a uniform $\delta_*>0$. Since
$B_{\kappa r/2}(x)\subset B_{\rho/2}(x)$, it follows that
\[
 M(\kappa r/2)\le(1-\delta_*)M(r).
\]
This comparison does not assume continuity at the frontier: the
harmonic replacement uses an actual Sobolev trace at an eligible
radius, whereas the indicator is used only in the Poisson
integral and is not asserted to belong to $H^{1/2}$.

Iteration, starting from $q_K\le1$, gives some
$\alpha\in(0,1)$ for which
\[
 \operatorname*{ess\,sup}_{B_r(x)}q_K
 \le C_*(r/r_*)^\alpha\qquad(0<r\le r_*).
\]
For example, decrease to an exponent smaller than one and no
larger than $\log(1-\delta_*)/\log(\kappa/2)$.
Interior smoothness now supplies a continuous representative
equal to zero at every frontier point. The strong maximum
principle gives $0<q_K<1$ on the connected exterior.
Uniqueness among continuous exterior solutions with these
boundary and end values follows by applying the maximum
principle on expanding truncated exteriors.

For completeness the same estimate gives a full spatial H\"older
modulus, not just decay at boundary points. If two points a
distance $\ell$ apart have one within $2\ell$ of $K$, both
values are bounded by $C_*\ell^\alpha$. Otherwise a local ball
joining them stays in the exterior. At distance $\varrho$ from
$K$, interior gradient estimates and the boundary decay give
$|\nabla q_K|\le C_*\varrho^{\alpha-1}$ when
$\varrho<r_*/4$; at larger distance fixed-radius interior
estimates suffice. Integration between the two points again
bounds their difference by $C_*\ell^\alpha$. Boundedness
handles larger separations. Hence the zero extensions satisfy
\begin{equation}
\label{eq:rpi:flow-holder-input}
 \|q_K\|_{C^{0,\alpha}(M,g_0)}\le C_*.
\end{equation}
The constants depend on $c$, the background geometry, and the initial
comparison radius. No derivatives of the weight enter this estimate.

\paragraph{The comparison for a H\"older weight.}
Suppose $\operatorname{osc}_{B_r}a\le C_a r^\alpha$. If
$P_{g_0}(L;B_r)\ge P_{g_0}(K;B_r)$ the conclusion is immediate.
Otherwise the density upper bound gives both perimeters at most
$C_*r^k$. The exact weighted comparison yields
\[
 \begin{aligned}
 P_{g_0}(K;B_r)
 &\le\frac{\sup_{B_r}a}{\inf_{B_r}a}P_{g_0}(L;B_r)\\
 &\le P_{g_0}(L;B_r)+C r^\alpha P_{g_0}(L;B_r)\\
 &\le P_{g_0}(L;B_r)+C r^{k+\alpha}.
 \end{aligned}
\]
This proves the final assertion.
\end{proof}

\subsection{The discrete flow and its area loss}

We now apply the static lemma successively. All enclosure minimizations
allow compact finite-perimeter supersets, including sets with bounded
complementary cavities, so union and intersection remain admissible.
Put $p=2k/d$, the exponent converting the conformal factor to its
perimeter weight.


Fix $0<\epsilon<1/2$, set $u_0^\epsilon=1$, and let
$K_0^\epsilon=K_0$. Inductively, put
\[
 h_j^\epsilon=(u_j^\epsilon)^{4/d}g_0,
 \qquad a_j^\epsilon=(1-\epsilon)^j.
\]
For $j\ge1$, let $K_j^\epsilon$ be the outermost
$h_j^\epsilon$-minimizing enclosure of $K_{j-1}^\epsilon$.
Let $v_j^\epsilon$ be the unique harmonic function on
$M\setminus K_j^\epsilon$ with zero boundary value and end value
$-a_j^\epsilon$, extended by zero to $K_j^\epsilon$. Define
\begin{equation}
\label{eq:rpi:flow-euler}
 \begin{aligned}
 u_{j+1}^\epsilon&=u_j^\epsilon+\epsilon v_j^\epsilon,\\
 u_t^\epsilon&=u_j^\epsilon+(t-j\epsilon)v_j^\epsilon,\\
 K_t^\epsilon&=K_j^\epsilon,\qquad
 v_t^\epsilon=v_j^\epsilon.
 \end{aligned}
\end{equation}
The last two lines apply when $j\epsilon\le t<(j+1)\epsilon$;
use the new hull and velocity at a grid endpoint.
Thus $u_t^\epsilon=1+\int_0^t v_s^\epsilon\,ds$ exactly.
The maximum principle gives
\[
 -a_j^\epsilon\le v_j^\epsilon\le0,
 \qquad
 (1-\epsilon)^{j+1}\le u_t^\epsilon\le1
 \quad(j\epsilon\le t<(j+1)\epsilon).
\]
More precisely, if $\theta=t-j\epsilon\in[0,\epsilon]$, then
concavity of $\log(1-\theta)$ gives
\[
 u_t^\epsilon\ge(1-\epsilon)^j(1-\theta)
 \ge(1-\epsilon)^{t/\epsilon}\ge4^{-t}.
\]
Consequently $u_t^\epsilon\ge4^{-T}$ for $0\le t\le T$.
In particular every finite interval has mesh-independent metric
ellipticity constants. The end values are exactly
\[
 u_j^\epsilon\longrightarrow a_j^\epsilon,
 \qquad
 u_t^\epsilon\longrightarrow
 a_j^\epsilon\bigl(1-(t-j\epsilon)\bigr).
\]
The discrete conductors are nested by construction. Comparison on
$M\setminus K_{j+1}^\epsilon$, including the change of the constant
at infinity, gives
\begin{equation}
\label{eq:rpi:flow-discrete-velocity-monotonicity}
 v_j^\epsilon\le v_{j+1}^\epsilon\le0.
\end{equation}
No limiting nesting is used here.


\begin{lemma}[Uniform estimates for the discrete flow]
\label{lem:rpi:discrete-flow}
For every finite $T$, the preceding induction defines the flow on
$[0,T]$. Every grid conductor $K_j^\epsilon$ is the outermost
$h_j^\epsilon$-minimizing enclosure of the fixed $K_0$. The conductors
are nested and have common compact support. Their velocities, and the
factors $u_t^\epsilon$, have a common spatial H\"older modulus;
$u_t^\epsilon$ is also uniformly Lipschitz in time. The frontiers have
uniform density and background almost-minimizing estimates.
Their grid areas $A_j^\epsilon=P_{h_j^\epsilon}(K_j^\epsilon)$
are nonincreasing, and the area is frozen between grid times.
\end{lemma}

\begin{proof}
\textbf{The fixed-obstacle minimizing property.}\quad
Every $K_j^\epsilon$ is also the outermost
$h_j^\epsilon$-minimizing enclosure of the fixed $K_0$.
Indeed, suppose this holds at step $j-1$, and let $L\supset K_0$
be any compact finite-perimeter competitor. The two consecutive metrics agree on
$K_{j-1}^\epsilon$, because $v_{j-1}^\epsilon$ vanishes there.
The boundary of $L\cap K_{j-1}^\epsilon$ is supported in that same
closed set. Hence
\[
 \begin{aligned}
 P_{h_j^\epsilon}(K_{j-1}^\epsilon)
 &=P_{h_{j-1}^\epsilon}(K_{j-1}^\epsilon)\\
 &\le P_{h_{j-1}^\epsilon}(L\cap K_{j-1}^\epsilon)
 =P_{h_j^\epsilon}(L\cap K_{j-1}^\epsilon).
 \end{aligned}
\]
Perimeter submodularity therefore gives
\[
 P_{h_j^\epsilon}(L\cup K_{j-1}^\epsilon)
 \le P_{h_j^\epsilon}(L).
\]
The set on the left is an admissible competitor for $K_j^\epsilon$,
which proves global minimality. If $L$ is any global minimizer,
the same inequalities show that its union with $K_{j-1}^\epsilon$
is a constrained minimizer. Outermostness of $K_j^\epsilon$ then
gives $L\subset K_j^\epsilon$. The induction starts at the actual
outermost initial hull, not at an arbitrary almost-minimizing
frontier.

Set $A_j^\epsilon=P_{h_j^\epsilon}(K_j^\epsilon)$. Freezing on the
previous conductor and its admissibility at the next step imply
\[
 0\le A_j^\epsilon\le A_{j-1}^\epsilon\le A_0.
\]
The metric at an intermediate time need not make the frozen
$K_j^\epsilon$ globally minimizing. Its area nevertheless equals
$A_j^\epsilon$, since $v_j^\epsilon=0$ on that frontier; comparison
with an arbitrary competitor always uses $h_j^\epsilon$.
On a fixed finite interval $h_j^\epsilon$ and $g_t^\epsilon$ differ
uniformly by $O_T(\epsilon)$ for
$j\epsilon\le t<(j+1)\epsilon$.


\paragraph{Closing the induction and its uniform estimates.}
At step $j$, the constructed continuous positive factor defines the
weighted hull problem. The fixed-obstacle minimizing property just
proved lets us apply Lemma~\ref{lem:rpi:static-hull-potential} with
$a=(u_j^\epsilon)^p$ and $c=4^{-pT}$. It gives the density estimates
and the continuous potential
\[
 v_j^\epsilon=-a_j^\epsilon q_{K_j^\epsilon},\qquad
 \|q_{K_j^\epsilon}\|_{C^{0,\alpha}}
       +\|v_j^\epsilon\|_{C^{0,\alpha}}\le C_T.
\]
In particular $v_j^\epsilon$ vanishes on the whole canonical conductor.
Thus the next factor is continuous and agrees with $u_j^\epsilon$
on $K_j^\epsilon$, including its perimeter traces. This closes the
induction underlying the freezing and fixed-obstacle comparisons.

The exact discrete integral identity gives a uniform spatial
H\"older bound for $u_t^\epsilon$ on $[0,T]$, while $|v_j^\epsilon|\le1$
gives the time Lipschitz bound. Hence
$\operatorname{osc}_{B_r}(u_j^\epsilon)^p\le C_T r^\alpha$.
The final assertion of the static lemma now gives the additive
almost-minimizing estimate with error $C_T r^{k+\alpha}$.

\paragraph{Uniform finite-time support.}
The global grid-minimality induction permits application of
\eqref{eq:rpi:flow-hull-radius} with the fixed obstacle $K_0$,
the fixed area bound $A_0$, and $c=4^{-pT}$. Hence, enlarging
the resulting radius if necessary,
\begin{equation}
\label{eq:rpi:flow-support}
 K_t^\epsilon\subset B_{R_T},\qquad
 \Vol_{g_0}(K_t^\epsilon)\le V_T
 \quad(0\le t\le T).
\end{equation}
It follows directly from the static truncation comparison and is
independent of the number or shape of the conductor components.

\end{proof}

\begin{proposition}[Vanishing discrete area loss]
\label{prop:rpi:discrete-area}
For every finite $T$, the areas of the discrete frontiers tend uniformly
to $A_0$ for $0\le t\le T$ as $\epsilon\downarrow0$.
\end{proposition}

\begin{proof}
Global minimality in the preceding grid metric, with
$K_j^\epsilon$ as competitor, gives
\[
 \begin{aligned}
 0\le A_{j-1}^\epsilon-A_j^\epsilon
 &\le\int_{\partial^*K_j^\epsilon}
       \bigl[(u_{j-1}^\epsilon)^p-(u_j^\epsilon)^p\bigr]
       \,d\mathcal H^k_{g_0}\\
 &=\int_{\partial^*K_j^\epsilon}
       \left[1-\left(1+
         \epsilon\frac{v_{j-1}^\epsilon}{u_{j-1}^\epsilon}
                    \right)^p\right]d\mathcal H^k_{h_{j-1}^\epsilon}\\
 &\le-p\epsilon\int_{\partial^*K_j^\epsilon}
          \frac{v_{j-1}^\epsilon}{u_{j-1}^\epsilon}
                   \,d\mathcal H^k_{h_{j-1}^\epsilon}\\
 &\le C_T\epsilon A_0
          \sup_{\partial K_j^\epsilon}|v_{j-1}^\epsilon|.
 \end{aligned}
\]
The third line follows from Bernoulli's inequality, with
$v_{j-1}^\epsilon\le0$.

Separate the grid transitions into those for which
\[
 \Vol_{g_0}(K_j^\epsilon\setminus K_{j-1}^\epsilon)
     >V_T\sqrt\epsilon
\]
and the remaining transitions. Discrete nesting and
\eqref{eq:rpi:flow-support} bound the number of the first kind by
$\epsilon^{-1/2}$, so their total area loss is $O_T(\sqrt\epsilon)$.
For a transition of the second kind, the two phase-density
inequalities imply
\[
 d_{\mathcal H}(\partial K_j^\epsilon,
                 \partial K_{j-1}^\epsilon)
 \le C_T\bigl(V_T\sqrt\epsilon\bigr)^{1/n}
\]
once $\epsilon$ is small. Indeed, a boundary point separated from
the other frontier by a ball of radius $\rho$ contributes at least
$c_T\rho^n$ to the swept volume: use the filled phase at a point of
the new frontier and the exterior phase at a point of the old one.
The old velocity vanishes on the old frontier; its H\"older modulus
therefore bounds its supremum on the new frontier by
$C_T\epsilon^{\alpha/(2n)}$. Summing the $\epsilon$-weighted
losses over at most $T/\epsilon+1$ transitions gives
\begin{equation}
\label{eq:rpi:flow-area-loss}
 0\le A_0-A_j^\epsilon
 \le C_T\bigl(\sqrt\epsilon+\epsilon^{\alpha/(2n)}\bigr)
 \longrightarrow0
\end{equation}
uniformly for $j\epsilon\le T$. The same assertion holds at
intermediate times because the current frontier has frozen area.
Only discrete nesting has entered this calculation.

\end{proof}

\subsection{The limiting minimum area, nesting, and integral evolution}

The discrete areas now converge to $A_0$. We first show that $A_0$ is
the minimum enclosure area in every limiting metric. This fact will
identify the selected outermost hulls and their velocities.

\begin{proof}[Proof of Lemma~\ref{lem:rpi:flow-construction}]
\textbf{Compactness and the limiting minimum area.}\quad
The functions $u_t^\epsilon$ are uniformly Lipschitz in time and
uniformly H\"older in space. Choose a mesh sequence tending to zero
for which they converge locally uniformly, jointly in space and
time, to a positive $u_t$, and put $g_t=u_t^{4/d}g_0$.
The support bound also controls the end normalization. Put
$\beta=d/2$ and enlarge $R_T$ until
$\Delta_{g_0}r^{-\beta}<0$ for $r\ge R_T$. The function
$1+v_t^\epsilon/a_j^\epsilon$ is harmonic there, lies between
zero and one on the inner sphere, and tends to zero at infinity.
The maximum principle gives
\[
 0\le v_j^\epsilon+a_j^\epsilon
       \le a_j^\epsilon(R_T/r)^\beta.
\]
Writing
$a_\epsilon(t)=a_j^\epsilon(1-t+j\epsilon)$ on the current
step and integrating yields
\[
 |u_t^\epsilon-a_\epsilon(t)|\le T(R_T/r)^\beta
 \qquad(0\le t\le T,\ r\ge R_T).
\]
The constants $a_\epsilon(t)$ converge uniformly to $e^{-t}$
on $[0,T]$. These tail bounds turn local uniform convergence of
the factors into uniform convergence on $[0,T]\times M$ and
retain the constant $-e^{-t}$ for every fixed-time limit of the
velocities. Thus the compactness passage preserves the end normalization.

We claim that, for every $t$, the minimum $g_t$-perimeter among
enclosures of $K_0$ equals $A_0$. Fix a time $t$ and a bounded
finite-perimeter enclosure $L\supset K_0$.
Grid minimality at $j=\lfloor t/\epsilon\rfloor$, the uniform
convergence of the lagged metric, and
\eqref{eq:rpi:flow-area-loss} imply
\begin{equation}
\label{eq:rpi:flow-competitor-lower-bound}
 P_{g_t}(L)\ge A_0.
\end{equation}
All discrete velocities vanish on $K_0$, so $u_t=1$ there and
$P_{g_t}(K_0)=A_0$. Hence $A_0$ is exactly the minimum enclosure
perimeter for $g_t$.

Common support and the background perimeter bound give a
subsequential BV limit $\widehat K_t$ of the grid conductors.
Lower semicontinuity in the uniformly converging positive weights
gives $P_{g_t}(\widehat K_t)\le A_0$. Since this is itself an
enclosure of $K_0$, \eqref{eq:rpi:flow-competitor-lower-bound}
gives the reverse inequality. Thus every such limit is a
$g_t$-minimizer of perimeter $A_0$. This also proves perimeter
convergence. The density estimates
then upgrade convergence of the canonical sets to Hausdorff
convergence of their frontiers.

Define $K_t$ to be the outermost $g_t$-minimizing enclosure of
$K_0$. It has perimeter $A_0$ and contains every $\widehat K_t$.
Lemma~\ref{lem:rpi:static-hull-potential} bounds its support by a
finite-time common radius and gives exact local weighted minimality,
since $u_t=1$ on $K_0$. It also supplies the uniform density, local
area, and background almost-minimizing estimates. A bounded component
of its exterior
could be filled with a strict decrease in perimeter; thus no such
component occurs. The unique asymptotically flat end accounts for
the remaining exterior component. These observations also apply
to every $\widehat K_t$.

\paragraph{Nesting of the limiting outermost hulls.}
For $s>t$, the monotonicity of the discrete conformal factors passes
to the limit, so $u_s\le u_t$. The already established minimum
areas give
\[
 A_0\le P_{g_s}(K_t)\le P_{g_t}(K_t)=A_0.
\]
It follows that $u_s=u_t$ almost everywhere on the perimeter of
$K_t$. Continuity of the factors and the positive perimeter density
at every frontier point improve this to
\begin{equation}
\label{eq:rpi:flow-frontier-freezing}
 u_s=u_t\quad\hbox{on }\partial K_t.
\end{equation}
At this stage no freezing throughout $K_t$ has been asserted.

Take any joint subsequential limit
$(\widehat K_s,\widehat v_s)$ of the grid conductors and velocities
at time $s$. The boundary H\"older estimate, interior harmonic
compactness, and the fixed-sphere tail comparison identify
$\widehat v_s$ as the harmonic function with zero extension to
$\widehat K_s$ and end value $-e^{-s}$. It is strictly negative
on its connected end exterior. On the other hand, the discrete
integral identity and
\eqref{eq:rpi:flow-discrete-velocity-monotonicity} give
\[
 u_t^\epsilon-u_s^\epsilon
   =-\int_t^s v_\tau^\epsilon\,d\tau
   \ge(s-t)(-v_s^\epsilon).
\]
Passing to this subsequence yields
\[
 u_t-u_s\ge(s-t)(-\widehat v_s)>0
       \quad\hbox{on }M\setminus\widehat K_s.
\]
Together with \eqref{eq:rpi:flow-frontier-freezing}, this implies
$\partial K_t\subset\widehat K_s$. The connected exterior of
$\widehat K_s$ agrees with the exterior of $K_t$ near infinity and
cannot cross $\partial K_t$, so it is contained in that exterior.
Therefore
\begin{equation}
\label{eq:rpi:flow-limit-squeeze}
 K_t\subset\widehat K_s\subset K_s\qquad(t<s).
\end{equation}
This proves nesting using only the discrete integral identity and
the previously proved area equality.

\paragraph{Identification of the integral flow.}
On each finite interval the function $V(t)=\Vol_{g_0}(K_t)$ is
bounded and nondecreasing. Its positive left jumps occur at
only countably many times. Let $\mathcal J$ be the union of
these exceptional sets over integer time intervals, together
with $\{0\}$. If $t\notin\mathcal J$, every fixed-time grid
limit satisfies, by \eqref{eq:rpi:flow-limit-squeeze},
\[
 \bigcup_{r<t}K_r\subset\widehat K_t\subset K_t.
\]
The union may be taken over rational $r<t$, and its volume is
the left limit of $V$. At a nonjump time this equals $V(t)$.
The canonical identity $K_t=\supp(\chi_{K_t}\,dV_{g_0})$
therefore gives
\[
 \overline{\bigcup_{r<t}K_r}=K_t:
\]
a point of $K_t$ outside that closure would have a neighborhood
of positive $K_t$-volume disjoint from the union, contradicting
equality of volumes. Since $\widehat K_t$ is closed, it equals
$K_t$.

Thus at every nonexceptional time all subsequential grid limits
agree. Compactness proves convergence of the entire chosen
sequence in $L^1$ and in Hausdorff distance of canonical sets
and frontiers. The latter assertion follows directly from the
two phase-density estimates: a boundary point remaining away
from the limiting boundary would have a ball wholly in one
limiting phase, contrary to the corresponding positive discrete
phase volume and $L^1$ convergence. The reverse direction uses
the two positive limiting phase volumes. Common support
prevents boundary points from escaping in this argument.
The boundary H\"older estimates, interior harmonic compactness,
uniqueness, and uniform tail comparison then give
\[
 \|v_t^\epsilon-v_t\|_{C^0(M)}\longrightarrow0
 \qquad(t\notin\mathcal J),
\]
where $v_t$ is the harmonic function on $E_t$, zero on $K_t$
and tending to $-e^{-t}$. The static density and capacitary
arguments construct this selected potential at exceptional
times as well. Nesting and the maximum principle make $v_t$
nondecreasing in time, so it is measurable. Dominated
convergence in the exact discrete integral identity gives
\[
 u_t=1+\int_0^t v_s\,ds.
\]
In fact the convergence of these integrals is uniform in space
and on $[0,T]$, since
\[
 \sup_{0\le t\le T}
 \left\|\int_0^t(v_s^\epsilon-v_s)\,ds\right\|_{C^0(M)}
 \le\int_0^T\|v_s^\epsilon-v_s\|_{C^0(M)}\,ds\longrightarrow0.
\]
The scalar norm is measurable by taking the supremum over a
countable dense spatial set; it is bounded by two, so the
displayed almost-everywhere convergence justifies the last
dominated-convergence step. The same exceptional set therefore works
for every spatial point.
The possible failure of grid velocities to converge to the selected
$v_t$ at a jump time has no effect on this integral. Now, and only
now, nesting and the integral identity give
$u_s=u_t$ throughout $K_t$ for $s\ge t$, since each later velocity
vanishes there.

The end normalization gives $v_s\ge-e^{-s}$, so the same
identity improves the finite-time lower bound to
\[
 u_t\ge1-\int_0^t e^{-s}\,ds=e^{-t}.
\]
It also shows that $u_t$ is $g_0$-harmonic in $E_t$:
each earlier velocity is harmonic there, and interior estimates
justify integration on compact subsets of that exterior. All
conductors contain the fixed collar of the physical boundary,
so $v_t=0$ and $u_t=1$ on that collar. No boundary condition or
mean-curvature assumption at the buried boundary is added.
The conformal transformation of the scalar Laplacian gives
\[
 \Delta_{g_t}\left(\frac{v_t}{u_t}\right)
 =u_t^{-(n+2)/d}
   \left(\Delta_{g_0}v_t-\frac{v_t}{u_t}\Delta_{g_0}u_t\right)=0
 \quad\hbox{in }E_t.
\]
Consequently $w_t=1+v_t/u_t$ is the $g_t$-harmonic capacitary
potential: its boundary value is one and its end value is zero.
The inequalities $-e^{-t}\le v_t\le0$ and $u_t\ge e^{-t}$
give $0\le w_t\le1$ directly. Its constant extension over $K_t$
equals one on the same buried collar.

Finally choose the convergent mesh subsequence diagonally over
$T=1,2,\ldots$. The true hull at each time is determined by its
limiting metric, so the constructions on overlapping finite
intervals agree. This gives the flow for all $t\ge0$.
Equation~\eqref{eq:rpi:flow-area} already holds at every time by
the minimizing-value argument, including jump times. The local
weighted comparison and its H\"older-weight conversion give
\eqref{eq:rpi:flow-almost-minimizing} uniformly on each finite
interval, as asserted.
\end{proof}
